\section{Applications: Quantization and Optimization}
\label{sec:applications}

We now demonstrate how averaging-based spike suppression can be used in practice
to address two major challenges induced by anisotropy in large language model
training: low-bit quantization and adaptive optimization.
In both settings, the dominant spectral spike is known to be harmful, and prior
solutions rely on explicit spectral operations.
Our analysis shows that simple averaging can fill the role of SVD-based spike
suppression, providing a lightweight and hardware-friendly alternative.

\subsection{Low-Bit Quantization}
\label{subsec:quantization}

\paragraph{Problem setting.}
In low-bit and mixed-precision training, quantization error is dominated by
outliers and large-magnitude directions.
When representations are anisotropic, the dominant spectral spike sets the
numerical scale for quantization, forcing the remaining long-tail directions
into a narrow dynamic range.
This results in severe precision loss for semantically meaningful components and
has been identified as a key bottleneck in FP4 and FP8 training.

Recent approaches such as Metis mitigate this issue by explicitly suppressing
spectral spikes using SVD-based projections.
While effective, these methods require expensive matrix decompositions and are
difficult to integrate efficiently into large-scale training pipelines.

\paragraph{Averaging as a quantization-friendly approximation.}
The analysis in Sections~3 and~4 shows that the batch mean $\mu$ lies almost
entirely in the spike subspace and that mean removal suppresses a coherent
direction responsible for a large fraction of spike energy.
At the same time, tail energy is preserved up to vanishing error.
This makes averaging a natural preconditioning step for quantization:
\begin{enumerate}
  \item subtract the batch mean to suppress the dominant spike;
  \item quantize the mean-removed representations;
  \item (optionally) add the mean back in higher precision if needed.
\end{enumerate}

Unlike SVD-based methods, averaging is rank-one, requires only reductions and
broadcasts, and maps efficiently to GPUs and accelerators.
Empirically, we observe that mean removal significantly reduces spike-dominated
scale effects while preserving long-tail resolution, closely mirroring the
behavior of spectral spike suppression.
This explains why averaging can serve as a practical surrogate for SVD-based
methods in low-bit training regimes.

\paragraph{Interpretation.}
From the perspective of Section~4, quantization is sensitive to the \emph{range}
of representations.
Averaging removes the shared mean direction that dominates this range, thereby
rebalancing dynamic range without collapsing semantic variation.
In this sense, averaging provides a minimal intervention that targets the root
cause of quantization instability: coherent anisotropy.

\subsection{Adaptive Optimization}
\label{subsec:optimization}

\paragraph{Problem setting.}
Adaptive optimizers such as Adam and AdamW normalize gradients by estimates of
second moments.
In the presence of a strong spectral spike, these estimates are dominated by
shared, low-rank components.
As a result, learning rates in long-tail directions are implicitly suppressed,
slowing semantic learning and reducing effective model capacity.

Recent spike-aware optimizers, including Spectra, explicitly identify and
attenuate dominant spectral components in gradient statistics.
While effective, these methods again rely on explicit spectral analysis or
structured approximations.

\paragraph{Averaging as spike-aware normalization.}
Our results show that the batch mean captures the dominant shared component in
both activations and gradients.
Subtracting this mean before forming second-moment statistics reduces the
influence of the spike on adaptive normalization:
\[
g_{\mathrm{avg}} = g - \mathbb{E}[g],
\]
where $g$ denotes per-token or per-batch gradients.
By Theorem~\ref{thm:avg_spectral}, this operation suppresses a coherent spike
direction while preserving the covariance of semantic residuals.

In effect, averaging performs a form of spike-aware normalization at negligible
cost.
It attenuates the shared component that would otherwise dominate second-moment
estimates, allowing adaptive optimizers to allocate learning rates more evenly
across semantic directions.

\paragraph{Relation to spike-aware optimizers.}
From a spectral perspective, averaging removes a rank-one approximation to the
spike subspace.
This is weaker than full spectral suppression but sufficient to address the most
harmful shared component.
Empirically, this conservative behavior aligns with the design philosophy of
Spectra: suppress dominant modes without amplifying noise or overcorrecting
tail directions.

Thus, averaging provides a simple mechanism that approximates the effect of
spike-aware optimizers, while retaining the stability and simplicity of standard
adaptive methods.

\subsection{Summary}

Across both quantization and optimization, anisotropy manifests as a coherent
spectral spike that degrades performance.
Our analysis shows that this spike can be isolated through simple averaging,
without explicit spectral decomposition.
Averaging therefore fills the role of SVD-based spike suppression in practice:
it is inexpensive, hardware-friendly, and preserves the semantic structure that
drives learning.

These applications illustrate the practical impact of the “blessing” aspect of
anisotropy: the same coherence that creates training pathologies also enables
efficient remedies when properly exploited.